\section{Cardinality of the distinct active directed rays}
\label{sec:area_law_initial_active_ray_cardinality}

Fix an arbitrary origin $o\in\mathbb R^2$, a finite endpoint set
$Z\subset\mathbb Z^2$, a natural width $C\ge2$, and an initial layer
$k_0\ge50{,}000{,}000$. Choose arbitrary valid pitch residues at every
layer. Let $\mathcal I$ be the resulting actual initial identifiers, and write
$X_i^0$ for their closed birth regions.
Let $k\ge k_0$, let $z\in F_{k,\ell_k}$ be an actual fine-cell index,
and let $v$ be one of the nine marks of its defining cell.
Put $\ell=\ell_k-5$ and $r=2^\ell/2=t_k/64$.

Use the all-midpoint fan of $v+[-r,r]^2$, with its eight-element
index set $\mathcal J$ and counterclockwise successor $\tau$.
Let $\sigma:\mathcal J\to\mathcal I$ be the unique actual assignment
of Theorem~\ref{thm:area_law_actual_sector_assignment}.
Write $P_s$ for the final endpoint of elementary base $s$, and write
$[P_s-v]_+$ for its directed ray. The endpoint vector is nonzero;
opposite vectors determine different directed rays. Define
\begin{align}
    A          & =\{s\in\mathcal J:\sigma(s)\ne\sigma(\tau(s))\}, \\
    \mathcal D & =\{[P_s-v]_+:s\in A\}.
\end{align}

% Source: 10-geometry.tex, geometry:initial-stars, lines 333--370,
% especially the opposite labels and single counting in lines 366--370;
% revision adc7f1241b42e322a6451854ab7e4b4c146bf78a.

\begin{theorem}[Even cardinality of the distinct active directed rays]
    \label{thm:area_law_initial_active_rays_cardinality}
    \lean{TNLean.PEPS.AreaLaw.Geometry.initialRegion_active_rays_card_even_and_le_eight}
    \leanok
    \uses{thm:area_law_actual_sector_assignment,
        def:area_law_cell_fan_successor,def:area_law_cell_fan_ray}
    The set of distinct directed rays at changes of the actual initial
    identifier has even cardinality at most eight:
    \begin{align}
        |\mathcal D| & \in2\mathbb N, & |\mathcal D| & \le8.
    \end{align}
\end{theorem}
\begin{proof}\leanok
    \uses{thm:area_law_cell_fan_ray_injective,
        thm:area_law_initial_sector_changes_even,def:area_law_cell_fan}
    Injectivity of the fan's directed rays gives $|\mathcal D|=|A|$.
    The parity theorem gives $|A|\in2\mathbb N$.
    Since $A\subseteq\mathcal J$ and $|\mathcal J|=8$,
    one also has $|\mathcal D|\le8$.
\end{proof}

\begin{theorem}[Actual frontier at changes of successive identifiers]
    \label{thm:area_law_initial_active_successor}
    \lean{TNLean.PEPS.AreaLaw.Geometry.initialRegion_frontier_near_mark_iff_active_successor}
    \leanok
    \uses{thm:area_law_actual_sector_assignment,
        def:area_law_cell_fan_successor}
    For every $x\in B_\infty(v,r)$, one has
    \begin{align}
        x\in\bigcup_{i\in\mathcal I}\partial X_i^0
         & \quad\Longleftrightarrow\quad
        \exists s\in A,\quad x\in[v,P_s].
    \end{align}
\end{theorem}
\begin{proof}\leanok
    \uses{thm:area_law_initial_active_radial,
        thm:area_law_cell_fan_successor_iff}
    The actual radial-frontier characterization gives two neighboring
    fan members $s,t$ whose common endpoint is $P_s$, whose assigned
    identifiers differ, and whose shared radial segment contains $x$.
    Endpoint adjacency is equivalent to $t=\tau(s)$.
    Thus these two members give precisely $s\in A$ and $x\in[v,P_s]$.
    Conversely, $s\in A$ and $x\in[v,P_s]$ give the neighboring pair
    $s,\tau(s)$ in the radial-frontier characterization.
\end{proof}
